
%==============================================================================
% CV Gabriel Reyna Alvarado
% ATS-friendly | Data Analyst / Machine Learning
% Compile with: pdflatex
%==============================================================================

\documentclass[10pt,letterpaper]{article}

%------------------------------------------------------------------------------
% PACKAGES
%------------------------------------------------------------------------------
\usepackage[T1]{fontenc}
\usepackage[utf8]{inputenc}
\usepackage{charter}
\usepackage[
    top=0.55in,
    bottom=0.55in,
    left=0.68in,
    right=0.68in
]{geometry}

\usepackage{titlesec}
\usepackage{enumitem}
\usepackage{xcolor}
\usepackage{hyperref}
\usepackage{microtype}

%------------------------------------------------------------------------------
% COLOURS
%------------------------------------------------------------------------------
\definecolor{ink}{HTML}{1A1A1A}
\definecolor{accent}{HTML}{14304A}
\definecolor{hairline}{HTML}{C8C3B8}
\definecolor{muted}{HTML}{555555}

%------------------------------------------------------------------------------
% PDF METADATA
%------------------------------------------------------------------------------
\hypersetup{
    colorlinks=true,
    urlcolor=accent,
    linkcolor=accent,
    pdftitle={Gabriel Reyna Alvarado - Curriculum Vitae},
    pdfauthor={Gabriel Reyna Alvarado},
    pdfsubject={Data Analyst / Machine Learning}
}

%------------------------------------------------------------------------------
% GLOBAL SETTINGS
%------------------------------------------------------------------------------
\AtBeginDocument{\color{ink}}
\pagestyle{empty}
\setlength{\parindent}{0pt}
\setlength{\parskip}{0pt}

%------------------------------------------------------------------------------
% SECTION HEADINGS
%------------------------------------------------------------------------------
\titleformat{\section}
    {\normalfont\normalsize\bfseries\scshape}
    {}{0pt}{}
    [\vspace{-0.65ex}{\color{hairline}\titlerule[0.45pt]}]

\titlespacing*{\section}
    {0pt}{1.65ex plus .2ex}{0.8ex}

%------------------------------------------------------------------------------
% ENTRY
%------------------------------------------------------------------------------
\newcommand{\entry}[2]{%
    \vspace{0.55ex}
    \textbf{#1}\par
    {\small\itshape\color{muted}#2}\par
    \vspace{0.15ex}
}

%------------------------------------------------------------------------------
% BULLETS
%------------------------------------------------------------------------------
\newenvironment{points}
{
    \begin{itemize}[
        leftmargin=1.25em,
        itemsep=0.18ex,
        parsep=0pt,
        topsep=0.25ex,
        label=\textbullet
    ]
}
{
    \end{itemize}
}

%------------------------------------------------------------------------------
% SKILL ROW
%------------------------------------------------------------------------------
\newcommand{\skillrow}[2]{%
    \textbf{#1:} #2\par
    \vspace{0.25ex}
}

%------------------------------------------------------------------------------
% C++
%------------------------------------------------------------------------------
% Charter deja un hueco visible entre los dos signos y sale como "C+ +".
% Kerning manual, sin paquetes extra. Si al compilar queda apretado o todavia
% suelto, ajustar estos dos valores.
\newcommand{\cpp}{C\kern-0.06em{+}\kern-0.10em{+}}

%------------------------------------------------------------------------------
% ENLACE CON DOMINIO VISIBLE
%------------------------------------------------------------------------------
% Un CV se imprime. Si la cabecera dice solo "Portfolio", en papel no lleva a
% ninguna parte. Esto mantiene el dominio a la vista y clicable a la vez.
% \urlref{https://ejemplo.com}{ejemplo.com}
\newcommand{\urlref}[2]{\href{#1}{#2}}

%==============================================================================
% DOCUMENT
%==============================================================================

\begin{document}

%------------------------------------------------------------------------------
% HEADER
%------------------------------------------------------------------------------

\begin{center}

    {\fontsize{20}{22}\selectfont\bfseries GABRIEL REYNA ALVARADO}

    \vspace{0.35ex}

    {\large\color{accent}
    Data Analyst \textbar{} Machine Learning
    }

    \vspace{0.65ex}

    {\small
    Lima, Peru
    \textbar{}
    +51 997 881 012
    \textbar{}
    \href{mailto:g.alonsoreyna@gmail.com}{g.alonsoreyna@gmail.com}
    }

    \vspace{0.3ex}

    {\small
    \urlref{https://portfolio-ztanq.vercel.app}{portfolio-ztanq.vercel.app}
    \textbar{}
    \urlref{https://linkedin.com/in/gabriel-reyna-alvarado}{linkedin.com/in/gabriel-reyna-alvarado}
    \textbar{}
    \urlref{https://github.com/ZtanQ}{github.com/ZtanQ}
    }

\end{center}

\vspace{0.35ex}

%------------------------------------------------------------------------------
% PROFILE
%------------------------------------------------------------------------------

\section{Profile}

Computer Science student with hands-on experience across the full data project
lifecycle: dataset preparation and cleaning, training and evaluation of
classification models, and deployment of results in dashboards for
non-technical users (Tableau and Streamlit). I work with explicit
methodological rigour --- data leakage detection, metric selection driven by
the cost of each error type, and documented limitations --- on a Python,
SQL and scikit-learn/XGBoost stack. Every project below links to its code and
to a write-up of the decisions behind it.

%------------------------------------------------------------------------------
% PROJECTS
%------------------------------------------------------------------------------

\section{Data and Machine Learning Projects}

\entry
{\href{https://portfolio-ztanq.vercel.app/projects/aldimi-predict}
{ALDIMI Predict --- Predictive System for a Paediatric Oncology Shelter}}
{Python, XGBoost, Scikit-learn, SHAP, SQLite, Streamlit
\textbar{} UPC, 2026
\textbar{} Team of four: modelling and deployment}

\begin{points}
    \item Trained and evaluated classification models (XGBoost, Random Forest)
    to anticipate 7- and 14-day stockouts in the inventory of a shelter for
    paediatric oncology patients, tuning hyperparameters with RandomizedSearchCV
    over 5-fold StratifiedKFold.

    \item Built a multiclass clinical risk classifier over 34 nutritional,
    socioeconomic and treatment-adherence variables, reaching a macro F1 of
    0.75 and AUC of 0.87 on a synthetic dataset generated to protect the
    confidentiality of real patients.

    \item Chose macro F1 over accuracy because the risk classes are imbalanced
    and a missed high-risk case costs far more than a false positive.

    \item Interpreted predictions with SHAP values to surface the dominant risk
    factors, enabling shelter staff to audit a model that informs care decisions.

    \item Deployed the system as a Streamlit dashboard over a unified SQLite
    database, delivered with a user manual and a data dictionary.
\end{points}


\entry
{\href{https://portfolio-ztanq.vercel.app/projects/smart-kitchen-intelligence}
{Smart Kitchen Intelligence --- Food Waste Analytics}}
{Python, Polars, Pandas, Scikit-learn, Tableau
\textbar{} UPC, 2026
\textbar{} Team of three: data preparation, modelling and visualisation}

\begin{points}
    \item Implemented the cleaning pipeline over 25,819 inventory movements,
    applying 6 verifiable quality rules (outliers, category harmonisation,
    referential integrity, structural nulls, deduplication), with
    cross-implementation in Polars and Pandas for independent validation.

    \item Trained waste classification models (Random Forest against Logistic
    Regression), detecting and removing a data leak that inflated metrics to
    F1 = 1.0, and selected the final model on F1 and cross-validation stability
    rather than accuracy, since a trivial classifier already reached 64.5\%.

    \item Designed a 10-sheet Tableau dashboard with 4 KPIs, a global household
    filter and a prescriptive panel, which identified that 97\% of the
    economic loss concentrated in a single storage location, with a waste rate
    9.1 times higher than the most efficient one.

    \item Validated the representativeness of the simulated dataset against
    the UNEP Food Waste Index Report 2024 benchmark, and documented the
    structural redundancy between location and product category as a
    methodological limitation.
\end{points}


\entry
{\href{https://portfolio-ztanq.vercel.app/projects/ski-recommender}
{Anti-Waste Recommendation Engine --- Hybrid Recommender}}
{Python, Scikit-learn, ALS, NetworkX, Streamlit
\textbar{} UPC, 2026
\textbar{} Team of three: modelling and deployment}

\begin{points}
    \item Trained a three-signal hybrid recommender combining TF-IDF content
    similarity, implicit ALS collaborative filtering, and expiry-proximity
    urgency, evaluated under a blind masked-basket completion protocol and
    achieving 100\% catalogue coverage against the popularity baseline.

    \item Segmented consumption patterns with DBSCAN after benchmarking against
    K-Means and GMM on the PCA-reduced space, reaching a Silhouette Score of
    0.65 with under 0.3\% noise.

    \item Empirically ruled out adding PageRank centrality to the ensemble
    through an ablation over the weight grid (optimal weight of zero),
    documented the negative result in the technical report, and deployed the
    system as a Streamlit application.
\end{points}


\entry
{\href{https://portfolio-ztanq.vercel.app/projects/fruitguard}
{FruitGuard --- Produce Freshness Classification with Computer Vision}}
{Python, OpenCV, Scikit-learn, Scikit-image, Gradio
\textbar{} University of Pittsburgh, 2025
\textbar{} Team of three: modelling and evaluation}

\begin{points}
    \item Built a freshness classifier for nine produce types over 30,357
    images, combining HOG, GLCM and colour features with an OpenCV preprocessing
    pipeline (normalisation and noise reduction), reaching 97.4\% test accuracy
    across 14 evaluated classes.
\end{points}

%------------------------------------------------------------------------------
% EXPERIENCE
%------------------------------------------------------------------------------

\section{Professional Experience}

\entry
{Data Analyst Intern --- Cirion Next Generation}
{Lima, Peru \textbar{} January 2024 -- March 2024}

\begin{points}
    \item Processed and cleaned datasets of over 10,000 records using SQL and
    Python (Pandas), implementing consistency checks that eliminated manual
    rework in engineering reports.

    \item Automated the generation of recurring visualisations with Matplotlib
    and Seaborn, freeing roughly 20 hours of manual work per month for the team.

    \item Refactored SQL queries alongside the data engineering team, reducing
    execution times for business-critical reports.
\end{points}


\entry
{Founder --- \href{https://www.camotestudio.com/}{Camote Studio}}
{Lima, Peru \textbar{} 2025 -- Present}

\begin{points}
    \item Co-founded an independent game studio and shipped \emph{Space Drunks},
    a cartoon-styled beat 'em up, working across development, production and
    team coordination rather than a single fixed role.

    \item Established the Git version control workflow for a repository of over
    500 binary assets, eliminating integration conflicts during team work.
\end{points}

%------------------------------------------------------------------------------
% TECHNICAL SKILLS
%------------------------------------------------------------------------------

\section{Technical Skills}

\skillrow{Languages}
{Python, SQL, Java, \cpp, R, JavaScript}

\skillrow{Data and Machine Learning}
{Pandas, Polars, NumPy, Scikit-learn, XGBoost, SHAP, OpenCV}

\skillrow{Techniques}
{Supervised classification, clustering (DBSCAN, K-Means, GMM), PCA, t-SNE,
collaborative filtering (ALS), TF-IDF, cross-validation}

\skillrow{Visualisation}
{Tableau, Streamlit, Matplotlib, Seaborn, Plotly}

\skillrow{Databases}
{SQL Server, MySQL, SQLite, MongoDB, relational modelling}

\skillrow{Tools}
{Git, GitHub, Docker, pytest, Excel}

\skillrow{Methodologies}
{CRISP-DM, Scrum}

%------------------------------------------------------------------------------
% EDUCATION
%------------------------------------------------------------------------------

\section{Education}

\entry
{BSc in Computer Science --- Universidad Peruana de Ciencias Aplicadas (UPC)}
{Lima, Peru \textbar{} Expected graduation: 2026}

Relevant coursework: Machine Learning, Data Visualisation, Database Management,
Statistical Analysis.


\entry
{Computer Vision \& AI Programme --- University of Pittsburgh}
{Remote \textbar{} 2022 -- 2025}

International collaboration programme. Where the FruitGuard classifier was built.


\entry
{Game Development Workshop --- UPC GameLab}
{Lima, Peru \textbar{} 2024 -- 2025}

Term-long workshop where teams form and ship a game; the team behind Camote
Studio came out of it. Modelled real-time simulation systems with
object-oriented programming and optimised data structures, reducing processing
overhead by 15\% against the baseline implementation.

%------------------------------------------------------------------------------
% LANGUAGES, CERTIFICATIONS AND ACTIVITIES
%------------------------------------------------------------------------------

\section{Languages, Certifications and Activities}

\skillrow{Languages}
{Spanish (native), English (B2 --- Cambridge First, 2025), French (intermediate)}

\skillrow{Certifications}
{Scrum Fundamentals Certified (SFC)
\textbar{}
\href{https://www.linkedin.com/in/gabriel-reyna-alvarado/details/certifications/}
{LinkedIn Certifications}}

\skillrow{Activities}
{PMI Lima Peru Chapter; PMI LATAM Conference 2026.
Inter-university innovation competition on SDG 11 (Sustainable Cities and
Communities): member of a team of 6 that developed and defended a proposal
before an industry jury.}

\end{document}
```
